\documentclass[10pt,a4paper]{amsart}
\usepackage{amsmath,amssymb,amsthm,graphicx}
% Retained minimal preamble for the Stacks Project chapter more-algebra.tex (More on Algebra).
% The chapter file pinned-original-source.tex carries only \input{preamble} on its first line, so no
% preamble text lives inside the chapter itself.  The Stacks Project uses one single shared preamble for
% every chapter; the snapshot stages it identically as sources/stacks-chow/preamble.tex and
% sources/stacks-etale-cohomology/preamble.tex.  The block below is copied mechanically, line for line,
% from sources/stacks-chow/preamble.tex lines 14-18 (the Xy-pic setup) and 154-217 (the native theorem
% declarations together with the mathematical macros).  Nothing was added, removed or redefined; the
% wrapper class and the amsmath/amssymb/amsthm/graphicx base packages are supplied by the assembler.
% Chapter: more-algebra.tex -- More on Algebra.
% For commutative diagrams we use Xy-pic
\usepackage[all]{xy}

% We use 2cell for 2-commutative diagrams.
\xyoption{2cell}

% Theorem environments.
%
\theoremstyle{plain}
\newtheorem{theorem}[subsection]{Theorem}
\newtheorem{proposition}[subsection]{Proposition}
\newtheorem{lemma}[subsection]{Lemma}

\theoremstyle{definition}
\newtheorem{definition}[subsection]{Definition}
\newtheorem{example}[subsection]{Example}
\newtheorem{exercise}[subsection]{Exercise}
\newtheorem{situation}[subsection]{Situation}

\theoremstyle{remark}
\newtheorem{remark}[subsection]{Remark}
\newtheorem{remarks}[subsection]{Remarks}

\numberwithin{equation}{subsection}

% Macros
%
\def\lim{\mathop{\mathrm{lim}}\nolimits}
\def\colim{\mathop{\mathrm{colim}}\nolimits}
\def\Spec{\mathop{\mathrm{Spec}}}
\def\Hom{\mathop{\mathrm{Hom}}\nolimits}
\def\Ext{\mathop{\mathrm{Ext}}\nolimits}
\def\SheafHom{\mathop{\mathcal{H}\!\mathit{om}}\nolimits}
\def\SheafExt{\mathop{\mathcal{E}\!\mathit{xt}}\nolimits}
\def\Sch{\mathit{Sch}}
\def\Mor{\mathop{\mathrm{Mor}}\nolimits}
\def\Ob{\mathop{\mathrm{Ob}}\nolimits}
\def\Sh{\mathop{\mathit{Sh}}\nolimits}
\def\NL{\mathop{N\!L}\nolimits}
\def\CH{\mathop{\mathrm{CH}}\nolimits}
\def\proetale{{pro\text{-}\acute{e}tale}}
\def\etale{{\acute{e}tale}}
\def\QCoh{\mathit{QCoh}}
\def\Ker{\mathop{\mathrm{Ker}}}
\def\Im{\mathop{\mathrm{Im}}}
\def\Coker{\mathop{\mathrm{Coker}}}
\def\Coim{\mathop{\mathrm{Coim}}}

% Boxtimes
%
\DeclareMathSymbol{\boxtimes}{\mathbin}{AMSa}{"02}

%
% Macros for moduli stacks/spaces
%
\def\QCohstack{\mathcal{QC}\!\mathit{oh}}
\def\Cohstack{\mathcal{C}\!\mathit{oh}}
\def\Spacesstack{\mathcal{S}\!\mathit{paces}}
\def\Quotfunctor{\mathrm{Quot}}
\def\Hilbfunctor{\mathrm{Hilb}}
\def\Curvesstack{\mathcal{C}\!\mathit{urves}}
\def\Polarizedstack{\mathcal{P}\!\mathit{olarized}}
\def\Complexesstack{\mathcal{C}\!\mathit{omplexes}}
% \Pic is the operator that assigns to X its picard group, usage \Pic(X)
% \Picardstack_{X/B} denotes the Picard stack of X over B
% \Picardfunctor_{X/B} denotes the Picard functor of X over B
\def\Pic{\mathop{\mathrm{Pic}}\nolimits}
\def\Picardstack{\mathcal{P}\!\mathit{ic}}
\def\Picardfunctor{\mathrm{Pic}}
\def\Deformationcategory{\mathcal{D}\!\mathit{ef}}

% --- appended 2026-10-12 (Chief): the project own definitions of the reference/slogan/history
% display environments, copied mechanically from the shared Stacks preamble lines 8-12
% (\usepackage{verbatim} plus the three \newenvironment lines, which the chapter preamble keeps
% verbatim from the project preamble). Nothing added, removed or redefined: these environments
% carry no mathematical content and the project itself typesets them as comments. Without these
% three lines a counted statement whose own text uses them fails with Environment slogan undefined.
\usepackage{verbatim}
\newenvironment{reference}{\comment}{\endcomment}
\newenvironment{slogan}{\comment}{\endcomment}
\newenvironment{history}{\comment}{\endcomment}
\usepackage{cleveref}
\crefname{originallemma}{original lemma}{original lemmas}
\Crefname{originallemma}{Original lemma}{Original lemmas}
\makeatletter
\newlabel{more-algebra-lemma-dumb-properties-henselization}{{07QM}{source}{}{sourcepointer-1}{}}
\newlabel{more-algebra-lemma-dumb-properties-henselization@cref}{{[originallemma][0][]07QM}{[1][0][]source}{}{}{}}
\makeatother
\crefname{originallemma}{original lemma}{original lemmas}
\Crefname{originallemma}{Original lemma}{Original lemmas}
\makeatletter
\newlabel{algebra-lemma-henselization}{{04GN}{source}{}{sourcepointer-2}{}}
\newlabel{algebra-lemma-henselization@cref}{{[originallemma][0][]04GN}{[1][0][]source}{}{}{}}
\makeatother
\crefname{originallemma}{original lemma}{original lemmas}
\Crefname{originallemma}{Original lemma}{Original lemmas}
\makeatletter
\newlabel{algebra-lemma-etale}{{00U2}{source}{}{sourcepointer-3}{}}
\newlabel{algebra-lemma-etale@cref}{{[originallemma][0][]00U2}{[1][0][]source}{}{}{}}
\makeatother
\crefname{originallemma}{original lemma}{original lemmas}
\Crefname{originallemma}{Original lemma}{Original lemmas}
\makeatletter
\newlabel{algebra-lemma-colimit-flat}{{05UT}{source}{}{sourcepointer-4}{}}
\newlabel{algebra-lemma-colimit-flat@cref}{{[originallemma][0][]05UT}{[1][0][]source}{}{}{}}
\makeatother
\crefname{originallemma}{original lemma}{original lemmas}
\Crefname{originallemma}{Original lemma}{Original lemmas}
\makeatletter
\newlabel{algebra-lemma-local-flat-ff}{{00HR}{source}{}{sourcepointer-5}{}}
\newlabel{algebra-lemma-local-flat-ff@cref}{{[originallemma][0][]00HR}{[1][0][]source}{}{}{}}
\makeatother
\crefname{originallemma}{original lemma}{original lemmas}
\Crefname{originallemma}{Original lemma}{Original lemmas}
\makeatletter
\newlabel{algebra-lemma-strict-henselization}{{04GP}{source}{}{sourcepointer-6}{}}
\newlabel{algebra-lemma-strict-henselization@cref}{{[originallemma][0][]04GP}{[1][0][]source}{}{}{}}
\makeatother
\crefname{originallemma}{original lemma}{original lemmas}
\Crefname{originallemma}{Original lemma}{Original lemmas}
\makeatletter
\newlabel{algebra-lemma-local-artinian-basis-when-flat}{{051F}{source}{}{sourcepointer-7}{}}
\newlabel{algebra-lemma-local-artinian-basis-when-flat@cref}{{[originallemma][0][]051F}{[1][0][]source}{}{}{}}
\makeatother
\crefname{originallemma}{original lemma}{original lemmas}
\Crefname{originallemma}{Original lemma}{Original lemmas}
\makeatletter
\newlabel{algebra-lemma-flat-base-change}{{00HI}{source}{}{sourcepointer-8}{}}
\newlabel{algebra-lemma-flat-base-change@cref}{{[originallemma][0][]00HI}{[1][0][]source}{}{}{}}
\makeatother
\title{Faithful flatness, maximal ideals and residue rings of henselization and strict henselization}
\author{The Stacks Project Authors}
\date{stacks-project:more-algebra.tex:local-snapshot-20261012; GNU FDL 1.2 or later}
\begin{document}
\maketitle
\setcounter{subsection}{1}
\noindent\textit{Editorial note.} All delivered mathematical text is the author's own, copied from the pinned Stacks Project chapter file more-algebra.tex; no statement or proof step was rewritten, reordered or removed.\par\medskip
\noindent\textit{Editorial note.} The chapter file pinned-original-source.tex carries only \textbackslash{}input\{preamble\} on its first line, so no preamble text lives inside the chapter. The retained file sources/stacks-more-algebra/retained-preamble.tex is a mechanical, line-for-line copy of the project's single shared preamble as staged in the snapshot (Xy-pic setup plus theorem declarations and mathematical macros). Nothing was added, removed or redefined, and the excerpt serves only as preamble, never as a span coordinate.\par\medskip
\noindent\textit{Editorial note.} Both counted spans start at the \textbackslash{}begin line of their environment, so the statement and the author's proof are delivered whole; surrounding prose of the section is deliberately outside the counted spans and is carried as standard prerequisites instead.\par\medskip
\noindent\textit{Editorial note.} The citation fields (source\_id, source\_version, authors, source\_url, license.name, source\_title) are identical to those of every other card this lane has produced from this chapter, so the batch selector treats them as one source.\par\medskip
\noindent\textit{Allowed standard prerequisites (editorial).}\begin{itemize}
\item The henselization is a colimit of etale algebras, and the strict henselization is a colimit of finite etale algebras. Algebra, Lemma \textbackslash{}ref\{algebra-lemma-henselization\} and Algebra, Lemma \textbackslash{}ref\{algebra-lemma-strict-henselization\}, cited at proof lines 10927-10929 and 10936-10937. Give the presentations from which flatness is read off.
\item Etale maps are flat, flatness is stable under colimits, and a flat local homomorphism is faithfully flat. Algebra, Lemma \textbackslash{}ref\{algebra-lemma-etale\}, Algebra, Lemma \textbackslash{}ref\{algebra-lemma-colimit-flat\}, and Algebra, Lemma \textbackslash{}ref\{algebra-lemma-local-flat-ff\}, cited at proof lines 10928-10934. Produce faithful flatness of the two ring maps.
\item A flat local ring map which is an isomorphism on residue fields makes the source a basis of the target as a module over each Artinian quotient. Algebra, Lemma \textbackslash{}ref\{algebra-lemma-local-artinian-basis-when-flat\}, cited at proof lines 10946-10948 and 10951-10954. Gives both the equality of the quotients and the freeness statement.
\item Base change of a flat ring map is flat. Algebra, Lemma \textbackslash{}ref\{algebra-lemma-flat-base-change\}, cited at proof lines 10945-10947. Makes the induced map on the quotients by powers of the maximal ideal flat.
\end{itemize}\medskip
\ifdefined\hypertarget\hypertarget{sourcepointer-1}{}\fi
\noindent\textit{Source reference.} Original lemma 07QM: focused-seeds/AGP/sources/stacks-snapshot-20261011/tags.txt, source lines 9837--9837. The lemma's own canonical tag, quoted so that readers can resolve the label inside the delivered text; it points at the counted statement itself and carries no proof content.\par\medskip
\ifdefined\hypertarget\hypertarget{sourcepointer-2}{}\fi
\noindent\textit{Source reference.} Original lemma 04GN: focused-seeds/AGP/sources/stacks-snapshot-20261011/tags.txt, source lines 5839--5839. Navigation pointer to a statement the counted proof refers to but does not contain; nothing load-bearing is delegated to it.\par\medskip
\ifdefined\hypertarget\hypertarget{sourcepointer-3}{}\fi
\noindent\textit{Source reference.} Original lemma 00U2: focused-seeds/AGP/sources/stacks-snapshot-20261011/tags.txt, source lines 1096--1096. Navigation pointer to a statement the counted proof refers to but does not contain; nothing load-bearing is delegated to it.\par\medskip
\ifdefined\hypertarget\hypertarget{sourcepointer-4}{}\fi
\noindent\textit{Source reference.} Original lemma 05UT: focused-seeds/AGP/sources/stacks-snapshot-20261011/tags.txt, source lines 7527--7527. Navigation pointer to a statement the counted proof refers to but does not contain; nothing load-bearing is delegated to it.\par\medskip
\ifdefined\hypertarget\hypertarget{sourcepointer-5}{}\fi
\noindent\textit{Source reference.} Original lemma 00HR: focused-seeds/AGP/sources/stacks-snapshot-20261011/tags.txt, source lines 653--653. Navigation pointer to a statement the counted proof refers to but does not contain; nothing load-bearing is delegated to it.\par\medskip
\ifdefined\hypertarget\hypertarget{sourcepointer-6}{}\fi
\noindent\textit{Source reference.} Original lemma 04GP: focused-seeds/AGP/sources/stacks-snapshot-20261011/tags.txt, source lines 5840--5840. Navigation pointer to a statement the counted proof refers to but does not contain; nothing load-bearing is delegated to it.\par\medskip
\ifdefined\hypertarget\hypertarget{sourcepointer-7}{}\fi
\noindent\textit{Source reference.} Original lemma 051F: focused-seeds/AGP/sources/stacks-snapshot-20261011/tags.txt, source lines 6533--6533. Navigation pointer to a statement the counted proof refers to but does not contain; nothing load-bearing is delegated to it.\par\medskip
\ifdefined\hypertarget\hypertarget{sourcepointer-8}{}\fi
\noindent\textit{Source reference.} Original lemma 00HI: focused-seeds/AGP/sources/stacks-snapshot-20261011/tags.txt, source lines 644--644. Navigation pointer to a statement the counted proof refers to but does not contain; nothing load-bearing is delegated to it.\par\medskip
\begin{lemma}
\label{lemma-dumb-properties-henselization}
Let $(R, \mathfrak m, \kappa)$ be a local ring. Then we have the following
\begin{enumerate}
\item $R \to R^h \to R^{sh}$ are faithfully flat ring maps,
\item $\mathfrak m R^h = \mathfrak m^h$ and
$\mathfrak m R^{sh} = \mathfrak m^h R^{sh} = \mathfrak m^{sh}$,
\item $R/\mathfrak m^n = R^h/\mathfrak m^nR^h$ for all $n$,
\item there exist elements $x_i \in R^{sh}$ such that
$R^{sh}/\mathfrak m^nR^{sh}$ is a free $R/\mathfrak m^n$-module
on $x_i \bmod \mathfrak m^nR^{sh}$.
\end{enumerate}
\end{lemma}

\begin{proof}
By construction $R^h$ is a colimit of \'etale $R$-algebras, see
Algebra, Lemma \ref{algebra-lemma-henselization}. Since \'etale
ring maps are flat (Algebra, Lemma \ref{algebra-lemma-etale}) we see that
$R^h$ is flat over $R$ by
Algebra, Lemma \ref{algebra-lemma-colimit-flat}.
As a flat local ring homomorphism is faithfully flat
(Algebra, Lemma \ref{algebra-lemma-local-flat-ff})
we see that $R \to R^h$ is faithfully flat.
The ring map $R^h \to R^{sh}$ is a colimit of finite \'etale ring
maps, see proof of Algebra, Lemma \ref{algebra-lemma-strict-henselization}.
Hence the same arguments as above show that $R^h \to R^{sh}$ is
faithfully flat.

\medskip\noindent
Part (2) follows from Algebra, Lemmas \ref{algebra-lemma-henselization} and
\ref{algebra-lemma-strict-henselization}. Part (3) follows from
Algebra, Lemma \ref{algebra-lemma-local-artinian-basis-when-flat}
because $R/\mathfrak m \to R^h/\mathfrak mR^h$ is an isomorphism and
$R/\mathfrak m^n \to R^h/\mathfrak m^nR^h$ is flat as a base change of
the flat ring map $R \to R^h$
(Algebra, Lemma \ref{algebra-lemma-flat-base-change}).
Let $\kappa^{sep}$ be the residue field of $R^{sh}$ (it is a
separable algebraic closure of $\kappa$). Choose $x_i \in R^{sh}$
mapping to a basis of $\kappa^{sep}$ as a $\kappa$-vector space.
Then (4) follows from
Algebra, Lemma \ref{algebra-lemma-local-artinian-basis-when-flat}
in exactly the same way as above.
\end{proof}

\end{document}
